\title{2020 年高考数学 全国II卷-理科}
\date{}
\maketitle

\section*{一、选择题}
本题共 12 小题，每小题 5 分，共 60 分.在每小题给出的四个选项中，只有一项是符合题目要求的.

\begin{question}[points=5]
已知集合 $U=\{-2,-1,0,1,2,3\}$，$A=\{-1,0,1\}$，$B=\{1,2\}$，则 $\complement_U(A\cup B)=$ （ ）
\begin{choices}
  \item $\{-2,3\}$
  \item $\{-2,2,3\}$
  \item $\{-2,-1,0,3\}$
  \item $\{-2,-1,0,2,3\}$
\end{choices}
\begin{solution}
\textbf{答案：}$\{-2,3\}$

\textbf{分析：}首先进行并集运算，然后计算补集即可.

\textbf{详解：}由题意可得：$A\cup B=\{-1,0,1,2\}$，则 $\complement_U(A\cup B)=\{-2,3\}$.
\end{solution}
\end{question}

\begin{question}[points=5]
若$\alpha$为第四象限角，则（ ）
\begin{choices}
  \item $\cos2\alpha>0$
  \item $\cos2\alpha<0$
  \item $\sin2\alpha>0$
  \item $\sin2\alpha<0$
\end{choices}
\begin{solution}
\textbf{答案：}$\sin2\alpha<0$

\textbf{分析：}由题意结合二倍角公式确定所给的选项是否正确即可.

\textbf{详解：}当 $\alpha=-\frac{\pi}{6}$ 时，$\cos2\alpha=\cos\left(-\frac{\pi}{3}\right)>0$，选项B错误；当 $\alpha=-\frac{\pi}{3}$ 时，$\cos2\alpha=\cos\left(-\frac{2\pi}{3}\right)<0$，选项A错误；由 $\alpha$ 在第四象限可得：$\sin\alpha<0,\cos\alpha>0$，则 $\sin2\alpha=2\sin\alpha\cos\alpha<0$，选项C错误，选项D正确.
\end{solution}
\end{question}

\begin{question}[points=5]
在新冠肺炎疫情防控期间，某超市开通网上销售业务，每天能完成1200份订单的配货，由于订单量大幅增加，导致订单积压.为解决困难，许多志愿者踊跃报名参加配货工作.已知该超市某日积压500份订单未配货，预计第二天的新订单超过1600份的概率为0.05，志愿者每人每天能完成50份订单的配货，为使第二天完成积压订单及当日订单的配货的概率不小于0.95，则至少需要志愿者（ ）
\begin{choices}
  \item 10名
  \item 18名
  \item 24名
  \item 32名
\end{choices}
\begin{solution}
\textbf{答案：}18名

\textbf{分析：}算出第二天订单数，除以志愿者每天能完成的订单配货数即可.

\textbf{详解：}由题意，第二天新增订单数为 $500+1600-1200=900$，故需要志愿者 $\dfrac{900}{50}=18$ 名.
\end{solution}
\end{question}

\begin{question}[points=5]
北京天坛的圜丘坛为古代祭天的场所，分上、中、下三层，上层中心有一块圆形石板(称为天心石)，环绕天心石砌9块扇面形石板构成第一环，向外每环依次增加9块，下一层的第一环比上一层的最后一环多9块，向外每环依次也增加9块，已知每层环数相同，且下层比中层多729块，则三层共有扇面形石板(不含天心石)（ ）
\begin{choices}
  \item 3699块
  \item 3474块
  \item 3402块
  \item 3339块
\end{choices}
\begin{solution}
\textbf{答案：}3402块

\textbf{分析：}第n环天石心块数为$a_n$，第一层共有n环，则$\{a_n\}$是以9为首项，9为公差的等差数列，设$S_n$为$\{a_n\}$的前n项和，由题意可得$S_{3n}-S_{2n}=S_{2n}-S_n+729$，解方程即可得到n，进一步得到$S_{3n}$.

\textbf{详解：}设第n环天石心块数为$a_n$，第一层共有n环，则$\{a_n\}$是以9为首项，9为公差的等差数列，$a_n=9+(n-1)\times9=9n$，设$S_n$为$\{a_n\}$的前n项和，则第一层、第二层、第三层的块数分别为$S_n,S_{2n}-S_n,S_{3n}-S_{2n}$，因为下层比中层多729块，所以$S_{3n}-S_{2n}=S_{2n}-S_n+729$，即$\frac{3n(9+27n)}{2}-\frac{2n(9+18n)}{2}=\frac{2n(9+18n)}{2}-\frac{n(9+9n)}{2}+729$，即$9n^2=729$，解得$n=9$，所以$S_{3n}=S_{27}=\frac{27(9+9\times27)}{2}=3402$.
\end{solution}
\end{question}

\begin{question}[points=5]
若过点$(2,1)$的圆与两坐标轴都相切，则圆心到直线$2x-y-3=0$的距离为（ ）
\begin{choices}
  \item $\dfrac{\sqrt5}{5}$
  \item $\dfrac{2\sqrt5}{5}$
  \item $\dfrac{3\sqrt5}{5}$
  \item $\dfrac{4\sqrt5}{5}$
\end{choices}
\begin{solution}
\textbf{答案：}$\dfrac{2\sqrt5}{5}$

\textbf{分析：}由题意可知圆心在第一象限，设圆心的坐标为$(a,a),a>0$，可得圆的半径为$a$，写出圆的标准方程，利用点$(2,1)$在圆上，求得实数$a$的值，利用点到直线的距离公式可求出圆心到直线$2x-y-3=0$的距离.

\textbf{详解：}由于圆上的点$(2,1)$在第一象限，若圆心不在第一象限，则圆与至少与一条坐标轴相交，不合乎题意，所以圆心必在第一象限，设圆心的坐标为$(a,a)$，则圆的半径为$a$，圆的标准方程为$(x-a)^2+(y-a)^2=a^2$.由题意可得$(2-a)^2+(1-a)^2=a^2$，即$a^2-6a+5=0$，解得$a=1$或$a=5$，所以圆心的坐标为$(1,1)$或$(5,5)$，圆心到直线$2x-y-3=0$的距离均为$d=\dfrac{|2\times a - a -3|}{\sqrt{2^2+(-1)^2}}=\dfrac{|a-3|}{\sqrt5}$，代入$a=1$或$a=5$得$d=\dfrac{2}{\sqrt5}=\dfrac{2\sqrt5}{5}$.
\end{solution}
\end{question}

\begin{question}[points=5]
数列$\{a_n\}$中，$a_1=2$，$a_{m+n}=a_m a_n$，若$a_{k+1}+a_{k+2}+\cdots+a_{k+10}=2^{15}-2^5$，则$k=$（ ）
\begin{choices}
  \item 2
  \item 3
  \item 4
  \item 5
\end{choices}
\begin{solution}
\textbf{答案：}4

\textbf{分析：}取$m=1$，可得出数列$\{a_n\}$是等比数列，求得数列$\{a_n\}$的通项公式，利用等比数列求和公式可得出关于k的等式，由$k\in\mathbb{N}^*$可求得k的值.

\textbf{详解：}在等式$a_{m+n}=a_m a_n$中，令$m=1$，可得$a_{n+1}=a_n a_1=2a_n$，$\therefore\dfrac{a_{n+1}}{a_n}=2$，所以数列$\{a_n\}$是以2为首项，以2为公比的等比数列，则$a_n=2\times2^{n-1}=2^n$，$\therefore a_{k+1}+a_{k+2}+\cdots+a_{k+10}=\dfrac{a_{k+1}(1-2^{10})}{1-2}=2^{k+1}(2^{10}-1)=2^5(2^{10}-1)$，$\therefore2^{k+1}=2^5$，则$k+1=5$，解得$k=4$.
\end{solution}
\end{question}

\begin{question}[points=5]
如图是一个多面体的三视图，这个多面体某条棱的一个端点在正视图中对应的点为$M$，在俯视图中对应的点为$N$，则该端点在侧视图中对应的点为（ ）

\begin{center}
  \includegraphics[width=.38\linewidth]{question-7-1.jpg}
\end{center}
\begin{choices}
  \item E
  \item F
  \item G
  \item H
\end{choices}
\begin{solution}
\textbf{答案：}E

\textbf{分析：}根据三视图，画出多面体立体图形，即可求得M点在侧视图中对应的点.

\textbf{详解：}根据三视图，画出多面体立体图形，图中标出了根据三视图M点所在位置，可知在侧视图中所对应的点为E.
\end{solution}
\end{question}

\begin{question}[points=5]
设$O$为坐标原点，直线$x=a$与双曲线$C:\dfrac{x^2}{a^2}-\dfrac{y^2}{b^2}=1(a>0,b>0)$的两条渐近线分别交于$D,E$两点，若$\triangle ODE$的面积为8，则$C$的焦距的最小值为（ ）
\begin{choices}
  \item 4
  \item 8
  \item 16
  \item 32
\end{choices}
\begin{solution}
\textbf{答案：}8

\textbf{分析：}因为$C:\dfrac{x^2}{a^2}-\dfrac{y^2}{b^2}=1(a>0,b>0)$，可得双曲线的渐近线方程是$y=\pm\dfrac{b}{a}x$，与直线$x=a$联立方程求得$D,E$两点坐标，即可求得$|ED|$，根据$\triangle ODE$的面积为8，可得$ab$值，根据$2c=2\sqrt{a^2+b^2}$，结合均值不等式，即可求得答案.

\textbf{详解：}双曲线的渐近线方程是$y=\pm\dfrac{b}{a}x$，直线$x=a$与渐近线交于$D(a,b),E(a,-b)$，所以$|ED|=2b$，$\triangle ODE$面积为$\dfrac{1}{2}a\times2b=ab=8$，双曲线焦距$2c=2\sqrt{a^2+b^2}\geq2\sqrt{2ab}=2\sqrt{16}=8$，当且仅当$a=b=2\sqrt2$时取等号，所以C的焦距的最小值为8.
\end{solution}
\end{question}

\begin{question}[points=5]
设函数$f(x)=\ln|2x+1|-\ln|2x-1|$，则$f(x)$（ ）
\begin{choices}
  \item 是偶函数，且在$(\frac12,+\infty)$单调递增
  \item 是奇函数，且在$(-\frac12,\frac12)$单调递减
  \item 是偶函数，且在$(-\infty,-\frac12)$单调递增
  \item 是奇函数，且在$(-\infty,-\frac12)$单调递减
\end{choices}
\begin{solution}
\textbf{答案：}是奇函数，且在$(-\infty,-\frac12)$单调递减

\textbf{分析：}根据奇偶性的定义可判断出$f(x)$为奇函数，排除AC；当$x\in(-\frac12,\frac12)$时，利用函数单调性的性质可判断出$f(x)$单调递增，排除B；当$x\in(-\infty,-\frac12)$时，利用复合函数单调性可判断出$f(x)$单调递减，从而得到结果.

\textbf{详解：}由$f(x)=\ln|2x+1|-\ln|2x-1|$得$f(x)$定义域为$\{x|x\neq\pm\frac12\}$，关于坐标原点对称，又$f(-x)=\ln|1-2x|-\ln|-2x-1|=\ln|2x-1|-\ln|2x+1|=-f(x)$，$\therefore f(x)$为定义域上的奇函数，排除AC；当$x\in(-\frac12,\frac12)$时，$f(x)=\ln(2x+1)-\ln(1-2x)$，$\because y=\ln(2x+1)$在$(-\frac12,\frac12)$上单调递增，$y=\ln(1-2x)$在$(-\frac12,\frac12)$上单调递减，$\therefore f(x)$在$(-\frac12,\frac12)$上单调递增，排除B；当$x\in(-\infty,-\frac12)$时，$f(x)=\ln(-2x-1)-\ln(1-2x)=\ln\dfrac{2x+1}{2x-1}=\ln\left(1+\dfrac{2}{2x-1}\right)$，$\because \mu=1+\dfrac{2}{2x-1}$在$(-\infty,-\frac12)$上单调递减，$f(\mu)=\ln\mu$在定义域内单调递增，根据复合函数单调性可知：$f(x)$在$(-\infty,-\frac12)$上单调递减，D正确.
\end{solution}
\end{question}

\begin{question}[points=5]
已知$\triangle ABC$是面积为$\dfrac{9\sqrt3}{4}$的等边三角形，且其顶点都在球$O$的球面上.若球$O$的表面积为$16\pi$，则$O$到平面$ABC$的距离为（ ）
\begin{choices}
  \item $\sqrt3$
  \item $\dfrac{3}{2}$
  \item 1
  \item $\dfrac{\sqrt3}{2}$
\end{choices}
\begin{solution}
\textbf{答案：}1

\textbf{分析：}根据球O的表面积和$\triangle ABC$的面积可求得球O的半径R和$\triangle ABC$外接圆半径r，由球的性质可知所求距离$d=\sqrt{R^2-r^2}$.

\textbf{详解：}设球O的半径为R，则$4\pi R^2=16\pi$，解得$R=2$.设$\triangle ABC$外接圆半径为r，边长为a，$\triangle ABC$是面积为$\dfrac{9\sqrt3}{4}$的等边三角形，$\therefore\dfrac12 a^2\times\dfrac{\sqrt3}{2}=\dfrac{9\sqrt3}{4}$，解得$a=3$，$\therefore r=\dfrac23\times\sqrt{a^2-\frac{a^2}{4}}=\dfrac23\times\sqrt{9-\frac94}=\dfrac23\times\frac{3\sqrt3}{2}=\sqrt3$，$\therefore$球心O到平面ABC的距离$d=\sqrt{R^2-r^2}=\sqrt{4-3}=1$.
\end{solution}
\end{question}

\begin{question}[points=5]
若$2^x-2^y<3^{-x}-3^{-y}$，则（ ）
\begin{choices}
  \item $\ln(y-x+1)>0$
  \item $\ln(y-x+1)<0$
  \item $\ln|x-y|>0$
  \item $\ln|x-y|<0$
\end{choices}
\begin{solution}
\textbf{答案：}$\ln(y-x+1)>0$

\textbf{分析：}将不等式变为$2^x-3^{-x}<2^y-3^{-y}$，根据$f(t)=2^t-3^{-t}$的单调性知$x<y$，以此去判断各个选项中真数与1的大小关系，进而得到结果.

\textbf{详解：}由$2^x-2^y<3^{-x}-3^{-y}$得：$2^x-3^{-x}<2^y-3^{-y}$，令$f(t)=2^t-3^{-t}$，$\because y=2^x$为R上的增函数，$y=3^{-x}$为R上的减函数，$\therefore f(t)$为R上的增函数，$\therefore x<y$，$\because y-x>0$，$\therefore y-x+1>1$，$\therefore\ln(y-x+1)>0$，则A正确，B错误；$\because|x-y|$与1的大小不确定，故CD无法确定.
\end{solution}
\end{question}

\begin{question}[points=5]
0-1周期序列在通信技术中有着重要应用.若序列$a_1a_2\cdots a_n\cdots$满足$a_i\in\{0,1\}(i=1,2,\cdots)$，且存在正整数$m$，使得$a_{i+m}=a_i(i=1,2,\cdots)$成立，则称其为0-1周期序列，并称满足$a_{i+m}=a_i(i=1,2,\cdots)$的最小正整数$m$为这个序列的周期.对于周期为$m$的0-1序列$a_1a_2\cdots a_n\cdots$，$C(k)=\dfrac1m\sum\limits_{i=1}^m a_i a_{i+k}(k=1,2,\cdots,m-1)$是描述其性质的重要指标，下列周期为5的0-1序列中，满足$C(k)\le\dfrac15(k=1,2,3,4)$的序列是（ ）
\begin{choices}
  \item $11010\cdots$
  \item $11011\cdots$
  \item $10001\cdots$
  \item $11001\cdots$
\end{choices}
\begin{solution}
\textbf{答案：}$10001\cdots$

\textbf{分析：}对于周期为5的序列，依次计算各选项的$C(1),C(2),C(3),C(4)$，判断是否都满足$C(k)\le\frac15$.

\textbf{详解：}由$a_{i+m}=a_i$知，序列$a_i$的周期为$m$，由已知$m=5$，$C(k)=\frac15\sum_{i=1}^5 a_i a_{i+k},k=1,2,3,4$. 对于选项A，$C(1)=\frac15(1+0+0+0+0)=\frac15\le\frac15$，$C(2)=\frac15(0+1+0+1+0)=\frac25>\frac15$，不满足；对于选项B，$C(1)=\frac15(1+0+0+1+1)=\frac35>\frac15$，不满足；对于选项C，$C(1)=\frac15(0+0+0+0+1)=\frac15$，$C(2)=\frac15(0+0+0+0+0)=0$，$C(3)=\frac15(0+0+0+0+1)=\frac15$，$C(4)=\frac15(1+0+0+0+0)=\frac15$，全部满足$\le\frac15$；对于选项D，$C(1)=\frac15(1+0+0+0+1)=\frac25>\frac15$，不满足. 故选C.
\end{solution}
\end{question}

\section*{二、填空题}
本题共4小题，每小题5分，共20分.

\begin{question}[points=5]
已知单位向量$\boldsymbol{a},\boldsymbol{b}$的夹角为$45^\circ$，$k\boldsymbol{a}-\boldsymbol{b}$与$\boldsymbol{a}$垂直，则$k=$.
\fillin[{$\dfrac{\sqrt2}{2}$}]
\begin{solution}
\textbf{答案：}$\dfrac{\sqrt2}{2}$

\textbf{分析：}首先求得向量的数量积，然后结合向量垂直的充分必要条件即可求得实数k的值.

\textbf{详解：}由题意可得：$\boldsymbol{a}\cdot\boldsymbol{b}=1\times1\times\cos45^\circ=\dfrac{\sqrt2}{2}$，由向量垂直的充分必要条件可得：$(k\boldsymbol{a}-\boldsymbol{b})\cdot\boldsymbol{a}=0$，即$k\boldsymbol{a}^2-\boldsymbol{a}\cdot\boldsymbol{b}=k-\dfrac{\sqrt2}{2}=0$，解得$k=\dfrac{\sqrt2}{2}$.
\end{solution}
\end{question}

\begin{question}[points=5]
4名同学到3个小区参加垃圾分类宣传活动，每名同学只去1个小区，每个小区至少安排1名同学，则不同的安排方法共有种.
\fillin[{$36$}]
\begin{solution}
\textbf{答案：}$36$

\textbf{分析：}根据题意，采用捆绑法，先取2名同学看作一组，现在可看成是3组同学分配到3个小区，即可求得答案.

\textbf{详解：}$\because$4名同学到3个小区参加垃圾分类宣传活动，每名同学只去1个小区，每个小区至少安排1名同学，$\therefore$先取2名同学看作一组，选法有$\mathrm{C}_4^2=6$，现在可看成是3组同学分配到3个小区，分法有$\mathrm{A}_3^3=6$，根据分步乘法原理，可得不同的安排方法$6\times6=36$种.
\end{solution}
\end{question}

\begin{question}[points=5]
设复数$z_1,z_2$满足$|z_1|=|z_2|=2$，$z_1+z_2=\sqrt3+\mathrm{i}$，则$|z_1-z_2|=$.
\fillin[{$2\sqrt3$}]
\begin{solution}
\textbf{答案：}$2\sqrt3$

\textbf{分析：}令$z_1=2\cos\theta+2\sin\theta\cdot\mathrm{i}$，$z_2=2\cos\alpha+2\sin\alpha\cdot\mathrm{i}$，根据复数相等可求得$\cos\theta\cos\alpha+\sin\theta\sin\alpha=-\frac12$，代入复数模长的公式中即可得到结果.

\textbf{详解：}$\because|z_1|=|z_2|=2$，可设$z_1=2\cos\theta+2\sin\theta\cdot\mathrm{i}$，$z_2=2\cos\alpha+2\sin\alpha\cdot\mathrm{i}$，$\therefore z_1+z_2=2(\cos\theta+\cos\alpha)+2(\sin\theta+\sin\alpha)\mathrm{i}=\sqrt3+\mathrm{i}$，$\therefore\begin{cases}2(\cos\theta+\cos\alpha)=\sqrt3\\2(\sin\theta+\sin\alpha)=1\end{cases}$，两式平方作和得：$4(2+2(\cos\theta\cos\alpha+\sin\theta\sin\alpha))=4$，化简得$\cos\theta\cos\alpha+\sin\theta\sin\alpha=-\frac12$，$\therefore|z_1-z_2|^2=4(\cos\theta-\cos\alpha)^2+4(\sin\theta-\sin\alpha)^2=8-8(\cos\theta\cos\alpha+\sin\theta\sin\alpha)=8+4=12$，$\therefore|z_1-z_2|=2\sqrt3$.
\end{solution}
\end{question}

\begin{question}[points=5]
设有下列四个命题：

$p_1$：两两相交且不过同一点的三条直线必在同一平面内.

$p_2$：过空间中任意三点有且仅有一个平面.

$p_3$：若空间两条直线不相交，则这两条直线平行.

$p_4$：若直线$l\subset$平面$\alpha$，直线$m\perp$平面$\alpha$，则$m\perp l$.

则下述命题中所有真命题的序号是.

①$p_1\land p_4$ ②$p_1\land p_2$ ③$\lnot p_2\lor p_3$ ④$\lnot p_3\lor\lnot p_4$
\fillin[{$①③④$}]
\begin{solution}
\textbf{答案：}$①③④$

\textbf{分析：}利用两交线直线确定一个平面可判断命题$p_1$的真假；利用三点共线可判断命题$p_2$的真假；利用异面直线可判断命题$p_3$的真假，利用线面垂直的定义可判断命题$p_4$的真假.再利用复合命题的真假可得出结论.

\textbf{详解：}对于命题$p_1$，可设$l_1$与$l_2$相交，这两条直线确定的平面为$\alpha$；若$l_3$与$l_1$相交，则交点$A$在平面$\alpha$内，同理，$l_3$与$l_2$的交点$B$也在平面$\alpha$内，所以$AB\subset\alpha$，即$l_3\subset\alpha$，命题$p_1$为真命题；对于命题$p_2$，若三点共线，则过这三个点的平面有无数个，命题$p_2$为假命题；对于命题$p_3$，空间中两条直线相交、平行或异面，命题$p_3$为假命题；对于命题$p_4$，若直线$m\perp$平面$\alpha$，则$m$垂直于平面$\alpha$内所有直线，$\because$直线$l\subset$平面$\alpha$，$\therefore$直线$m\perp$直线$l$，命题$p_4$为真命题.综上可知，$p_1\land p_4$为真命题，$p_1\land p_2$为假命题，$\lnot p_2\lor p_3$为真命题，$\lnot p_3\lor\lnot p_4$为真命题.故答案为：①③④.
\end{solution}
\end{question}

\section*{三、解答题}
共70分.解答应写出文字说明、证明过程或演算步骤.第17～21题为必考题，每个试题考生都必须作答.第22、23题为选考题，考生根据要求作答.

\begin{question}[points=12]
$\triangle ABC$中，$\sin^2A-\sin^2B-\sin^2C=\sin B\sin C$.

（1）求$A$；

（2）若$BC=3$，求$\triangle ABC$周长的最大值.
\begin{solution}
\textbf{答案：}（1）$\dfrac{2\pi}{3}$；（2）$3+2\sqrt3$

\textbf{分析：}（1）利用正弦定理角化边，配凑出$\cos A$的形式，进而求得$A$；（2）利用余弦定理可得到$(AC+AB)^2-AC\cdot AB=9$，利用基本不等式可求得$AC+AB$的最大值，进而得到结果.

\textbf{详解：}（1）由正弦定理可得：$BC^2-AC^2-AB^2=AC\cdot AB$，$\therefore\cos A=\dfrac{AC^2+AB^2-BC^2}{2AC\cdot AB}=-\dfrac12$，$\because A\in(0,\pi)$，$\therefore A=\dfrac{2\pi}{3}$.
（2）由余弦定理得：$BC^2=AC^2+AB^2-2AC\cdot AB\cos A=AC^2+AB^2+AC\cdot AB=9$，即$(AC+AB)^2-AC\cdot AB=9$，$\because AC\cdot AB\leq\left(\dfrac{AC+AB}{2}\right)^2$（当且仅当$AC=AB$时取等号），$\therefore9=(AC+AB)^2-AC\cdot AB\geq(AC+AB)^2-\left(\dfrac{AC+AB}{2}\right)^2=\dfrac34(AC+AB)^2$，解得：$AC+AB\leq2\sqrt3$（当且仅当$AC=AB$时取等号），$\therefore\triangle ABC$周长$L=AC+AB+BC\leq3+2\sqrt3$，$\therefore\triangle ABC$周长的最大值为$3+2\sqrt3$.
\end{solution}
\end{question}

\begin{question}[points=12]
某沙漠地区经过治理，生态系统得到很大改善，野生动物数量有所增加.为调查该地区某种野生动物的数量，将其分成面积相近的200个地块，从这些地块中用简单随机抽样的方法抽取20个作为样区，调查得到样本数据$(x_i,y_i)(i=1,2,\ldots,20)$，其中$x_i$和$y_i$分别表示第$i$个样区的植物覆盖面积(单位：公顷)和这种野生动物的数量，并计算得$\sum\limits_{i=1}^{20}x_i=60$，$\sum\limits_{i=1}^{20}y_i=1200$，$\sum\limits_{i=1}^{20}(x_i-\overline{x})^2=80$，$\sum\limits_{i=1}^{20}(y_i-\overline{y})^2=9000$，$\sum\limits_{i=1}^{20}(x_i-\overline{x})(y_i-\overline{y})=800$.

（1）求该地区这种野生动物数量的估计值（这种野生动物数量的估计值等于样区这种野生动物数量的平均数乘以地块数）；

（2）求样本$(x_i,y_i)(i=1,2,\ldots,20)$的相关系数（精确到0.01）；

（3）根据现有统计资料，各地块间植物覆盖面积差异很大.为提高样本的代表性以获得该地区这种野生动物数量更准确的估计，请给出一种你认为更合理的抽样方法，并说明理由.

附：相关系数$r=\dfrac{\sum\limits_{i=1}^{n}(x_i-\overline{x})(y_i-\overline{y})}{\sqrt{\sum\limits_{i=1}^{n}(x_i-\overline{x})^2\sum\limits_{i=1}^{n}(y_i-\overline{y})^2}}$，$\sqrt2=1.414$.
\begin{solution}
\textbf{答案：}（1）12000；（2）0.94；（3）分层抽样，理由见解析.

\textbf{分析：}（1）利用野生动物数量的估计值等于样区野生动物平均数乘以地块数，代入数据即可；（2）利用公式$r=\dfrac{\sum(x_i-\overline{x})(y_i-\overline{y})}{\sqrt{\sum(x_i-\overline{x})^2\sum(y_i-\overline{y})^2}}$计算即可；（3）各地块间植物覆盖面积差异较大，为提高样本数据的代表性，应采用分层抽样.

\textbf{详解：}（1）样区野生动物平均数为$\dfrac1{20}\sum_{i=1}^{20}y_i=\dfrac{1}{20}\times1200=60$，地块数为200，该地区这种野生动物数量的估计值为$200\times60=12000$.
（2）样本$(x_i,y_i)$的相关系数为$r=\dfrac{\sum_{i=1}^{20}(x_i-\overline{x})(y_i-\overline{y})}{\sqrt{\sum_{i=1}^{20}(x_i-\overline{x})^2\sum_{i=1}^{20}(y_i-\overline{y})^2}}=\dfrac{800}{\sqrt{80\times9000}}=\dfrac{800}{\sqrt{720000}}=\dfrac{800}{600\sqrt2}=\dfrac{4}{3\sqrt2}=\dfrac{2\sqrt2}{3}\approx0.94$.
（3）由于各地块间植物覆盖面积差异较大，为提高样本的代表性，应采用分层抽样.先将植物覆盖面积按优中差分成三层，在各层内按比例抽取样本，在每层内用简单随机抽样法抽取样本即可.
\end{solution}
\end{question}

\begin{question}[points=12]
已知椭圆$C_1:\dfrac{x^2}{a^2}+\dfrac{y^2}{b^2}=1(a>b>0)$的右焦点$F$与抛物线$C_2$的焦点重合，$C_1$的中心与$C_2$的顶点重合.过$F$且与$x$轴垂直的直线交$C_1$于$A,B$两点，交$C_2$于$C,D$两点，且$|CD|=\dfrac43|AB|$.

（1）求$C_1$的离心率；

（2）设$M$是$C_1$与$C_2$的公共点，若$|MF|=5$，求$C_1$与$C_2$的标准方程.
\begin{solution}
\textbf{答案：}（1）$\dfrac12$；（2）$C_1:\dfrac{x^2}{36}+\dfrac{y^2}{27}=1$，$C_2:y^2=12x$

\textbf{分析：}（1）求出$AB$、$CD$，利用$CD=\frac43AB$可得出关于$a,c$的齐次等式，可解得椭圆$C_1$的离心率的值；（2）由（1）可得出$C_1$的方程为$\dfrac{x^2}{4c^2}+\dfrac{y^2}{3c^2}=1$，联立曲线$C_1$与$C_2$的方程，求出点$M$的坐标，利用抛物线的定义结合$|MF|=5$可求得$c$的值，进而可得出$C_1$与$C_2$的标准方程.

\textbf{详解：}（1）$\because F(c,0)$，$AB\perp x$轴且与椭圆$C_1$相交于$A,B$两点，则直线$AB$的方程为$x=c$，联立$\begin{cases}x=c\\\dfrac{x^2}{a^2}+\dfrac{y^2}{b^2}=1\end{cases}$，解得$y=\pm\dfrac{b^2}{a}$，则$|AB|=\dfrac{2b^2}{a}$，抛物线$C_2$的方程为$y^2=4cx$，联立$\begin{cases}x=c\\y^2=4cx\end{cases}$，解得$y=\pm2c$，$\therefore |CD|=4c$，$\because CD=\dfrac43AB$，即$4c=\dfrac43\times\dfrac{2b^2}{a}$，$\therefore 2b^2=3ac$，即$2(a^2-c^2)=3ac$，即$2c^2+3ac-2a^2=0$，即$2e^2+3e-2=0$，$\because 0<e<1$，解得$e=\dfrac12$.
（2）由（1）知$a=2c$，$b=\sqrt3c$，椭圆$C_1$的方程为$\dfrac{x^2}{4c^2}+\dfrac{y^2}{3c^2}=1$，联立$\begin{cases}y^2=4cx\\\dfrac{x^2}{4c^2}+\dfrac{y^2}{3c^2}=1\end{cases}$，消去$y$并整理得$3x^2+16cx-12c^2=0$，解得$x=\dfrac23c$或$x=-6c$（舍去），由抛物线的定义可得$|MF|=\dfrac23c+c=\dfrac{5c}{3}=5$，解得$c=3$，$\therefore$曲线$C_1$的标准方程为$\dfrac{x^2}{36}+\dfrac{y^2}{27}=1$，曲线$C_2$的标准方程为$y^2=12x$.
\end{solution}
\end{question}

\begin{question}[points=12]
如图，已知三棱柱$ABC-A_1B_1C_1$的底面是正三角形，侧面$BB_1C_1C$是矩形，$M,N$分别为$BC,B_1C_1$的中点，$P$为$AM$上一点，过$B_1C_1$和$P$的平面交$AB$于$E$，交$AC$于$F$.

（1）证明：$AA_1\parallel MN$，且平面$A_1AMN\perp$平面$EB_1C_1F$；

（2）设$O$为$\triangle A_1B_1C_1$的中心，若$AO\parallel$平面$EB_1C_1F$，且$AO=AB$，求直线$B_1E$与平面$A_1AMN$所成角的正弦值.
\begin{solution}
\textbf{答案：}（1）证明见解析；（2）$\dfrac{\sqrt{10}}{10}$

\textbf{分析：}（1）由$M,N$分别为$BC,B_1C_1$的中点，$MN\parallel CC_1$，根据条件可得$AA_1\parallel BB_1$，可证$MN\parallel AA_1$，要证平面$EB_1C_1F\perp$平面$A_1AMN$，只需证明$EF\perp$平面$A_1AMN$即可；（2）连接$NP$，先求证四边形$ONPA$是平行四边形，根据几何关系求得$EP$，在$B_1C_1$截取$B_1Q=EP$，由（1）$BC\perp$平面$A_1AMN$，可得$\angle QPN$为$B_1E$与平面$A_1AMN$所成角，即可求得答案.

\textbf{详解：}（1）$\because M,N$分别为$BC,B_1C_1$的中点，$\therefore MN\parallel BB_1$，又$AA_1\parallel BB_1$，$\therefore MN\parallel AA_1$，在$\triangle ABC$中，$M$为$BC$中点，则$BC\perp AM$，又$\because$侧面$BB_1C_1C$为矩形，$\therefore BC\perp BB_1$，$\because MN\parallel BB_1$，$\therefore MN\perp BC$，由$MN\cap AM=M$，$MN,AM\subset$平面$A_1AMN$，$\therefore BC\perp$平面$A_1AMN$，又$\because B_1C_1\parallel BC$，且$B_1C_1\not\subset$平面$ABC$，$BC\subset$平面$ABC$，$\therefore B_1C_1\parallel$平面$ABC$，又$\because B_1C_1\subset$平面$EB_1C_1F$，且平面$EB_1C_1F\cap$平面$ABC=EF$，$\therefore B_1C_1\parallel EF$，$\therefore EF\parallel BC$，又$\because BC\perp$平面$A_1AMN$，$\therefore EF\perp$平面$A_1AMN$，$\because EF\subset$平面$EB_1C_1F$，$\therefore$平面$EB_1C_1F\perp$平面$A_1AMN$.
（2）连接$NP$，$\because AO\parallel$平面$EB_1C_1F$，平面$AONP\cap$平面$EB_1C_1F=NP$，$\therefore AO\parallel NP$，根据三棱柱上下底面平行，其面$A_1NMA\cap$平面$ABC=AM$，面$A_1NMA\cap$平面$A_1B_1C_1=A_1N$，$\therefore ON\parallel AP$，故四边形$ONPA$是平行四边形.设$\triangle ABC$边长是$6m(m>0)$，可得：$ON=AP$，$NP=AO=AB=6m$，$\because O$为$\triangle A_1B_1C_1$的中心，且$\triangle A_1B_1C_1$边长为$6m$，$\therefore ON=\frac13\times6\times\sin60^\circ=3m$，故$ON=AP=3m$，$\because EF\parallel BC$，$\therefore\dfrac{AP}{AM}=\dfrac{EP}{BM}$，$\therefore\dfrac{3m}{3\sqrt3m}=\dfrac{EP}{3m}$，解得$EP=m$，在$B_1C_1$截取$B_1Q=EP=m$，故$QN=2m$，$\because B_1Q=EP$且$B_1Q\parallel EP$，$\therefore$四边形$B_1QPE$是平行四边形，$\therefore B_1E\parallel PQ$，由（1）$B_1C_1\perp$平面$A_1AMN$，故$\angle QPN$为$B_1E$与平面$A_1AMN$所成角，在$\text{Rt}\triangle QPN$中，根据勾股定理可得：$PQ=\sqrt{QN^2+PN^2}=\sqrt{(2m)^2+(6m)^2}=2\sqrt{10}m$，$\therefore\sin\angle QPN=\dfrac{QN}{PQ}=\dfrac{2m}{2\sqrt{10}m}=\dfrac{\sqrt{10}}{10}$，$\therefore$直线$B_1E$与平面$A_1AMN$所成角的正弦值为$\dfrac{\sqrt{10}}{10}$.
\end{solution}
\end{question}

\begin{question}[points=12]
已知函数$f(x)=\sin^2x\sin2x$.

（1）讨论$f(x)$在区间$(0,\pi)$的单调性；

（2）证明：$|f(x)|\leq\dfrac{3\sqrt3}{8}$；

（3）设$n\in\mathbb{N}^*$，证明：$\sin^2x\sin^22x\sin^24x\cdots\sin^22^nx\leq\dfrac{3^n}{4^n}$.
\begin{solution}
\textbf{答案：}（1）当$x\in(0,\frac\pi3)$时，$f'(x)>0,f(x)$单调递增，当$x\in(\frac\pi3,\frac{2\pi}3)$时，$f'(x)<0,f(x)$单调递减，当$x\in(\frac{2\pi}3,\pi)$时，$f'(x)>0,f(x)$单调递增.（2）证明见解析；（3）证明见解析.

\textbf{分析：}(1)首先求得导函数的解析式，然后由导函数的零点确定其在各个区间上的符号，最后确定原函数的单调性即可；(2)首先确定函数的周期性，然后结合(1)中的结论确定函数在一个周期内的最大值和最小值即可证得题中的不等式；(3)对所给的不等式左侧进行恒等变形，然后结合(2)的结论和三角函数的有界性进行放缩即可证得题中的不等式.

\textbf{详解：}（1）$f(x)=\sin^2x\sin2x=2\sin^3x\cos x$，则$f'(x)=2(3\sin^2x\cos^2x-\sin^4x)=2\sin^2x(3\cos^2x-\sin^2x)=2\sin^2x(4\cos^2x-1)=2\sin^2x(2\cos x+1)(2\cos x-1)$，$f'(x)=0$在$x\in(0,\pi)$上的根为$x_1=\frac\pi3,x_2=\frac{2\pi}3$，当$x\in(0,\frac\pi3)$时，$f'(x)>0,f(x)$单调递增，当$x\in(\frac\pi3,\frac{2\pi}3)$时，$f'(x)<0,f(x)$单调递减，当$x\in(\frac{2\pi}3,\pi)$时，$f'(x)>0,f(x)$单调递增.
（2）注意到$f(x+\pi)=\sin^2(x+\pi)\sin[2(x+\pi)]=\sin^2x\sin2x=f(x)$，故函数$f(x)$是周期为$\pi$的函数，结合(1)的结论，计算可得：$f(0)=f(\pi)=0$，$f(\frac\pi3)=(\frac{\sqrt3}{2})^2\times\frac{\sqrt3}{2}=\frac{3\sqrt3}{8}$，$f(\frac{2\pi}3)=(\frac{\sqrt3}{2})^2\times(-\frac{\sqrt3}{2})=-\frac{3\sqrt3}{8}$，据此可得：$[f(x)]_{\max}=\frac{3\sqrt3}{8}$，$[f(x)]_{\min}=-\frac{3\sqrt3}{8}$，即$|f(x)|\leq\frac{3\sqrt3}{8}$.
（3）$\sin^2x\sin^22x\sin^24x\cdots\sin^22^nx=\left[\sin^3x\sin^32x\sin^34x\cdots\sin^32^nx\right]^{\frac23}=\left[\sin x(\sin^2x\sin2x)(\sin^22x\sin4x)\cdots(\sin^22^{n-1}x\sin2^nx)\sin^22^nx\right]^{\frac23}\leq\left[|\sin x|\times\frac{3\sqrt3}{8}\times\frac{3\sqrt3}{8}\times\cdots\times\frac{3\sqrt3}{8}\times|\sin2^nx|\right]^{\frac23}=\left[\left(\frac{3\sqrt3}{8}\right)^n\right]^{\frac23}=\left(\frac{3}{4}\right)^n$.
\end{solution}
\end{question}

\begin{question}[points=10]
已知曲线$C_1,C_2$的参数方程分别为$C_1:\begin{cases}x=4\cos^2\theta\\y=4\sin^2\theta\end{cases}$（$\theta$为参数），$C_2:\begin{cases}x=t+\dfrac1t\\y=t-\dfrac1t\end{cases}$（$t$为参数）.

（1）将$C_1,C_2$的参数方程化为普通方程；

（2）以坐标原点为极点，$x$轴正半轴为极轴建立极坐标系.设$C_1,C_2$的交点为$P$，求圆心在极轴上，且经过极点和$P$的圆的极坐标方程.
\begin{solution}
\textbf{答案：}（1）$C_1:x+y=4$；$C_2:x^2-y^2=4$；（2）$\rho=\dfrac{17}{5}\cos\theta$

\textbf{分析：}（1）分别消去参数$\theta$和$t$即可得到所求普通方程；（2）两方程联立求得点$P$，求得所求圆的直角坐标方程后，根据直角坐标与极坐标的互化即可得到所求极坐标方程.

\textbf{详解：}（1）由$\cos^2\theta+\sin^2\theta=1$得$C_1$的普通方程为$x+y=4$；由$\begin{cases}x=t+\frac1t\\y=t-\frac1t\end{cases}$得：$\begin{cases}x^2=t^2+\frac1{t^2}+2\\y^2=t^2+\frac1{t^2}-2\end{cases}$，两式相减得$C_2$的普通方程为$x^2-y^2=4$.
（2）由$\begin{cases}x+y=4\\x^2-y^2=4\end{cases}$得：$\begin{cases}x=\frac52\\y=\frac32\end{cases}$，即$P(\frac52,\frac32)$；设所求圆圆心的直角坐标为$(a,0)$，其中$a>0$，则$(a-\frac52)^2+(0-\frac32)^2=a^2$，解得$a=\frac{17}{10}$，$\therefore$所求圆的半径$r=a=\frac{17}{10}$，$\therefore$所求圆的直角坐标方程为$(x-\frac{17}{10})^2+y^2=(\frac{17}{10})^2$，即$x^2+y^2=\frac{17}{5}x$，$\therefore$所求圆的极坐标方程为$\rho=\frac{17}{5}\cos\theta$.
\end{solution}
\end{question}

\begin{question}[points=10]
已知函数$f(x)=|x-a^2|+|x-2a+1|$.

（1）当$a=2$时，求不等式$f(x)\ge4$的解集；

（2）若$f(x)\ge4$，求$a$的取值范围.
\begin{solution}
\textbf{答案：}（1）$\{x|x\le\dfrac32\text{ 或 }x\ge\dfrac{11}{2}\}$；（2）$(-\infty,-1]\cup[3,+\infty)$

\textbf{分析：}（1）分别在$x\le3$、$3<x<4$和$x\ge4$三种情况下解不等式求得结果；（2）利用绝对值三角不等式可得到$f(x)\ge(a-1)^2$，由此构造不等式求得结果.

\textbf{详解：}（1）当$a=2$时，$f(x)=|x-4|+|x-3|$.当$x\le3$时，$f(x)=4-x+3-x=7-2x\ge4$，解得$x\le\dfrac32$；当$3<x<4$时，$f(x)=4-x+x-3=1\ge4$，无解；当$x\ge4$时，$f(x)=x-4+x-3=2x-7\ge4$，解得$x\ge\dfrac{11}{2}$；综上所述：$f(x)\ge4$的解集为$\{x|x\le\frac32\text{ 或 }x\ge\frac{11}{2}\}$.
（2）$f(x)=|x-a^2|+|x-2a+1|\geq|(x-a^2)-(x-2a+1)|=|-(a^2-2a+1)|=(a-1)^2$（当且仅当$(x-a^2)(x-2a+1)\leq0$时取等号），$\therefore(a-1)^2\ge4$，解得：$a\le-1$或$a\ge3$，$\therefore a$的取值范围为$(-\infty,-1]\cup[3,+\infty)$.
\end{solution}
\end{question}

